% =============================================================================
% per_process_results.tex
%
% Per-process GAN performance summary and discussion of cross-process
% inconsistency. Standalone fragment intended to be \input{} into Report.tex.
% =============================================================================

\section{Per-Process Performance}
\label{sec:per-process}

The GAN was trained independently on $10^5$ Monte Carlo simulations of the
2-inverter circuit for each of the nine PTM process models. All runs used the
same architecture, learning-rate schedule, and 100-epoch budget. Final
metrics are summarized in Table~\ref{tab:per-process-metrics}. In every run
the training loss and average test loss are within $\sim$1\% of each other,
indicating the model is not overfitting --- the cross-process variation is
driven by representation and target-distribution effects rather than
generalization gap.

\begin{table}[h]
\centering
\caption{Per-process GAN performance after 100 epochs on $10^5$ simulations
of the 2-inverter circuit. Rows are sorted by mean relative error.}
\label{tab:per-process-metrics}
\begin{tabular}{lccc}
\toprule
Process       & Train Loss & Test Loss & Mean Rel. Err. ($I_{\text{target}}$) \\
\midrule
22nm\_LP      & 0.157 & 0.153 & 0.75  \\
32nm\_LP      & 0.118 & 0.120 & 1.27  \\
22nm\_HP      & 0.221 & 0.220 & 1.84  \\
45nm\_LP      & 0.116 & 0.116 & 2.64  \\
65nm\_bulk    & 0.178 & 0.180 & 2.93  \\
32nm\_HP      & 0.187 & 0.189 & 3.22  \\
45nm\_HP      & 0.202 & 0.201 & 9.06  \\
90nm\_bulk    & 0.236 & 0.236 & 9.72  \\
130nm\_bulk   & 0.306 & 0.306 & 36.72 \\
\bottomrule
\end{tabular}
\end{table}

\subsection{Why Performance Varies Across Processes}
\label{sec:why-variance}

Although the loss values cluster within a factor of two, the mean relative
error spans nearly two orders of magnitude. Several factors contribute,
roughly in decreasing order of suspected impact.

\paragraph{Target dynamic range differs by process.}
The PTM models span $V_{\text{th0}}$ from roughly $0.38\,\text{V}$ at 130nm
to $0.69\,\text{V}$ at 22nm\_LP, with corresponding differences in $V_{dd}$,
mobility $\mu_0$, and saturation velocity $V_{sat}$. Consequently
$I_{\text{target}}$ spans very different magnitudes across processes. Mean
\emph{relative} error is dominated by samples whose true current is small,
so a fixed absolute residual produces wildly different MRE values across
processes. The fact that MSE-style losses agree within $2\times$ while MRE
varies by $\sim 50\times$ is a strong indicator that the headline metric, not
the model, is producing most of the apparent variance.

\paragraph{Variable input dimensionality.}
Each PTM \texttt{.pm} file holds a different subset of BSIM parameters
constant. The dataset finalization step removes constant columns, so the
encoder sees a different feature-vector width per process. Hyperparameters
(\texttt{embedding\_dim}, \texttt{hidden\_dim}, attention heads) were tuned
for 22nm\_LP, which may not be appropriate for processes whose surviving
feature set is meaningfully different.

\paragraph{Physics-regime differences.}
HP versus LP variants differ substantially in short-channel and leakage
behavior. Older bulk nodes (65nm, 90nm, 130nm) operate in a long-channel
saturated regime where the current depends on parameters that the GAT may
weigh differently after per-process feature standardization.

\paragraph{Per-process label normalization is decoupled from the loss.}
The training pipeline fits \texttt{label\_log\_mean} and
\texttt{label\_log\_std} independently for each dataset, so MSE is taken in
a per-process log-z-score space. When predictions are back-transformed to
compute MRE, processes with wider log-$\sigma$ amplify an identical
z-score residual into a much larger relative error.

\paragraph{Fixed training budget for unequal difficulty.}
Every run used the same 100 epochs and learning-rate schedule. Training
losses for the older bulk processes were still trending down at the final
LR step (epoch 97), suggesting they are budget-limited rather than
converged.

\subsection{Proposed Improvements}
\label{sec:future-improvements}

The following directions are proposed for future work to reduce
cross-process variance and improve overall accuracy.

\begin{itemize}
  \item \textbf{Predict normalized current.} Replace the raw
    $I_{\text{target}}$ target with a process-normalized quantity such as
    $I / (W \cdot \mu C_{\text{ox}} \cdot V_{dd}^2)$ or $I / I_{\text{on,ref}}$
    so target distributions are comparable across processes.
  \item \textbf{Use a scale-aware loss.} Replace MSE on log-z-scored current
    with log-MAE or SMAPE so the loss directly tracks the relative-error
    metric reported in the table.
  \item \textbf{Standardize the input feature set.} Keep the union of all
    BSIM parameters (zero-padding where constant per process) so the
    encoder always sees a fixed-width input. This also enables
    cross-process training.
  \item \textbf{Process-conditioned training on the union of datasets.}
    Concatenate all nine datasets and condition the model on a one-hot or
    learned process embedding. With $\sim 9 \times 10^5$ samples the model
    can learn process-invariant physics while specializing per process via
    the conditioning vector. This is the most structurally promising
    change.
  \item \textbf{Shared encoder, per-process heads.} An alternative to
    conditioning: a shared graph encoder feeding small, per-process
    regression heads. Captures shared physics with process-specific output
    scaling.
  \item \textbf{Longer training and cosine schedule.} The LR was still
    halving and the loss still decreasing at the end of the 100-epoch
    runs; a cosine schedule with a longer budget should help especially on
    older bulk nodes.
  \item \textbf{Inspect tail behavior.} Report median and 90th-percentile
    relative error in addition to the mean to determine how much of the
    variance is driven by a small number of low-current outliers.
  \item \textbf{asinh transform for near-zero currents.} An inverse
    hyperbolic sine transform of the target handles small or negative
    currents gracefully and avoids the divergent behavior of $\log$ near
    zero.
\end{itemize}
